%======================================================================
%  LLM inference foundations
%======================================================================

\section{From text to autoregressive inference}
\label{sec:llm-inference-foundations}

A large language model is a neural sequence model with many learned numerical
parameters that assigns probabilities to token sequences.  ``Large'' has no
fixed parameter threshold; it indicates that model capacity and its associated
compute and memory demands have been scaled substantially.  This thesis treats
those already learned parameters as fixed inputs to the inference runtime;
how they were estimated is outside the scope of the work.  The question here
is what happens after a user supplies a prompt, which objects must be held in
memory, and which bytes the processor must move to produce each output token.

\subsection{Tokens, probabilities and generation}

A language model does not consume words directly.  A tokenizer maps text to
a sequence of integer identifiers drawn from a finite vocabulary
\(\mathcal{V}\).  A token can be a word, a word fragment, punctuation or even
whitespace, depending on the tokenizer.  Consequently, token count rather
than word count is the relevant unit for context capacity and inference
work.  Given a prefix \(x_{<t}=(x_1,\ldots,x_{t-1})\), an autoregressive
language model assigns a probability to every possible next token:

\begin{equation}
\label{eq:autoregressive-factorization}
  p_\theta(x_{1:T})
  = \prod_{t=1}^{T} p_\theta(x_t \mid x_{<t}),
  \qquad x_t \in \mathcal{V}.
\end{equation}
\eqcaption{Autoregressive factorization.  Generation is a sequence of
next-token decisions: every selected token becomes part of the prefix for the
following decision.}

The model computes one unnormalised score, or \emph{logit}, for every token in
\(\mathcal{V}\).  A softmax converts the logits into a probability
distribution, and a decoding policy selects a token.  Greedy decoding takes
the most probable token; sampling policies may adjust the distribution with
temperature, top-\(k\), or top-\(p\) filtering before drawing from it.  The
selected token is appended to the context and the model is evaluated again.
The loop ends at an end-of-sequence token, a requested length, or another
stopping rule.  Detokenization finally maps the generated identifiers back to
text.  \figref{fig:token-generation} shows this feedback loop.  The important
systems fact is that a response is not produced in one pass: its tokens are
serially dependent, even though the operations inside one pass are parallel.

\begin{figure}[htbp]
\centering
\resizebox{\textwidth}{!}{%
\begin{tikzpicture}[
  font=\small,
  ticket/.style={draw=black!55, fill=white, rounded corners=1.5pt,
                 minimum height=0.72cm, inner xsep=7pt, align=center},
  generated/.style={ticket, draw=hbmcol!85!black, fill=hbmcol!15},
  model/.style={draw=ddrcol!85!black, fill=ddrcol!12,
                rounded corners=3pt, minimum width=3.0cm,
                minimum height=1.0cm, align=center},
  flow/.style={-{Stealth}, thick, black!65},
  loop/.style={-{Stealth}, thick, hbmcol!85!black},
  note/.style={font=\footnotesize, text=black!72, align=center}
]
  \node[note, anchor=east] at (-6.35,1.2) {prompt tokens};
  \node[ticket] (p1) at (-5.55,1.2) {High};
  \node[ticket, right=0.05cm of p1] (p2) {bandwidth};
  \node[ticket, right=0.05cm of p2] (p3) {memory};
  \node[ticket, right=0.05cm of p3] (p4) {can};

  \node[model] (model) at (0.1,-0.35)
    {decoder-only\\transformer};
  \node[model, right=1.0cm of model] (logits)
    {logits over\\the vocabulary};
  \node[model, right=1.0cm of logits] (select)
    {decoding policy\\select one token};
  \node[generated, right=0.9cm of select] (new) {reduce};

  \draw[flow] (p4.south) |- (model.north);
  \draw[flow] (model) -- (logits);
  \draw[flow] (logits) -- (select);
  \draw[flow] (select) -- (new);
  \draw[loop] (new.south) -- ++(0,-0.85) -| (model.south)
    node[pos=0.28, below, note, text=hbmcol!85!black]
      {append to the context; repeat};

  \node[note, anchor=east] at (-6.35,-2.45) {next context};
  \node[ticket] at (-5.55,-2.45) {High};
  \node[ticket] at (-3.75,-2.45) {bandwidth};
  \node[ticket] at (-1.95,-2.45) {memory};
  \node[ticket] at (-0.65,-2.45) {can};
  \node[generated] at (0.65,-2.45) {reduce};
  \node[note, anchor=west, text width=4.9cm] at (1.55,-2.45)
    {The new token is now an input to the next step.};
\end{tikzpicture}%
}
\caption{Autoregressive generation at token level.  Token boundaries are
illustrative and depend on the tokenizer; the example is not intended to show
the output of a particular model.  Author's synthesis.}
\label{fig:token-generation}
\end{figure}

The \emph{context} at a given step comprises the prompt tokens and all tokens
generated so far.  Its maximum permitted length is the context window.  The
prompt length and requested generation length are separate workload
variables: the former controls how much history must initially be processed,
whereas the latter controls how many serial decode iterations are executed.
Batch size or concurrency determines how many sequences share an evaluation
step.  These quantities must not be collapsed into a model name when reporting
performance.

\subsection{What one model evaluation computes}

Most contemporary generative LLMs use a decoder-only transformer derived from
the attention architecture of Vaswani et al.~\cite{vaswani2017attention}.
Token identifiers are mapped to vectors, position information is incorporated,
and the vectors pass through a stack of transformer layers.  Each layer
contains causal self-attention, a feed-forward network and the associated
normalisation and residual operations.  The causal constraint prevents a
position from using later tokens.

For one attention head, learned projections form queries, keys and values,
\(Q=XW_Q\), \(K=XW_K\), and \(V=XW_V\).  A simplified attention operation is

\begin{equation}
\label{eq:scaled-dot-product-attention}
  \operatorname{Attention}(Q,K,V)
  = \operatorname{softmax}\!\left(
      \frac{QK^{\mathsf T}}{\sqrt{d_h}} + M_{\mathrm{causal}}
    \right)V,
\end{equation}
\eqcaption{Scaled dot-product causal attention.  The mask excludes future
positions; during decode, the new query uses keys and values retained from the
existing context.}

where \(d_h\) is the head dimension and \(M_{\mathrm{causal}}\) represents the
causal mask.  Multi-head attention performs several such projections, and the
feed-forward block applies large learned linear transformations independently
at each token position.  After the final layer, an output projection produces
the vocabulary logits used by the decoding policy.  Architectural variants
change details such as positional encoding, activation function, normalisation
and the number of key--value heads, but not the data-lifetime distinction used
in this thesis.

That distinction has three parts~\cite{pope2023efficient,sheng2023flexgen}:

\begin{itemize}
  \item \textbf{Weights} are the fixed learned parameters.  They persist for
        the model's lifetime, are reused by all requests and usually dominate
        its capacity footprint.
  \item The \textbf{key--value (KV) cache} retains the key and value vectors
        already computed in every layer.  It avoids recomputing the complete
        prefix at each decode step, but is private to a sequence, grows with
        context length and concurrency, and must be read when later tokens
        attend to their history.
  \item \textbf{Activations and workspaces} are temporary intermediates and
        runtime buffers.  Their sizes depend on the phase, batch, kernels and
        implementation; prepacked or duplicated weights may add another
        runtime-managed allocation.
\end{itemize}

The KV cache is therefore not the model's memory.  It is request state beside
the model, and the two objects have different lifetimes and access patterns.
PagedAttention demonstrates how variable-size KV allocations and
fragmentation can limit serving capacity even when the arithmetic is
unchanged~\cite{kwon2023pagedattention}.  That study targets GPU serving, so it
supports the memory-management mechanism rather than a quantitative claim
about the CPU studied here.

\subsection{The inference pipeline and its memory objects}

\figref{fig:inference-memory-lifetimes} follows one request through the
runtime.  Model loading and any weight repacking happen before steady-state
generation.  Tokenization produces the prompt identifiers.  \emph{Prefill}
processes those identifiers and creates the initial KV cache.  The
\emph{decode} loop then evaluates one new position per sequence, selects a
token and appends its key and value state.  Sampling and detokenization are
usually small compared with the transformer layers, but they are inside the
user-visible latency boundary unless the measurement deliberately excludes
them.

\begin{figure}[htbp]
\centering
\resizebox{\textwidth}{!}{%
\begin{tikzpicture}[
  font=\small,
  stage/.style={draw=black!60, fill=neutralcol, rounded corners=2pt,
                minimum width=2.45cm, minimum height=0.95cm, align=center},
  flow/.style={-{Stealth}, thick, black!65},
  state/.style={rounded corners=2pt, minimum height=0.54cm, align=center,
                font=\footnotesize},
  label/.style={font=\footnotesize\bfseries, anchor=east},
  note/.style={font=\scriptsize, align=center, text=black!72}
]
  \node[stage] (load) at (-5.6,2.2) {load / initialise};
  \node[stage] (tok) at (-2.75,2.2) {tokenize prompt};
  \node[stage] (pre) at (0.1,2.2) {prefill};
  \node[stage] (dec) at (2.95,2.2) {decode loop};
  \node[stage] (out) at (5.8,2.2) {detokenize};
  \draw[flow] (load) -- (tok);
  \draw[flow] (tok) -- (pre);
  \draw[flow] (pre) -- (dec);
  \draw[flow] (dec) -- (out);
  \draw[flow] (dec.north) .. controls +(0,0.7) and +(0,0.7) .. (dec.north)
    node[above, note] {one token at a time};

  \node[label] at (-6.9,0.85) {weights};
  \node[state, draw=ddrcol!85!black, fill=ddrcol!13,
        minimum width=11.35cm, anchor=west] at (-5.9,0.85)
    {persistent, read-only model state reused by every request and decode step};

  \node[label] at (-6.9,0.0) {KV cache};
  \node[state, draw=hbmcol!85!black, fill=hbmcol!15,
        minimum width=2.45cm, anchor=west] (kvpre) at (-1.2,0.0)
    {created in prefill};
  \node[state, draw=hbmcol!85!black, fill=hbmcol!24,
        minimum width=5.45cm, anchor=west] at (1.4,0.0)
    {read and extended as the context grows};

  \node[label] at (-6.9,-0.85) {temporary};
  \node[state, draw=black!55, fill=black!6,
        minimum width=2.45cm, anchor=west] at (-4.05,-0.85)
    {token buffers};
  \node[state, draw=black!55, fill=black!6,
        minimum width=2.45cm, anchor=west] at (-1.2,-0.85)
    {prefill activations};
  \node[state, draw=black!55, fill=black!6,
        minimum width=2.45cm, anchor=west] at (1.65,-0.85)
    {step activations};
  \node[state, draw=black!55, fill=black!6,
        minimum width=2.45cm, anchor=west] at (4.5,-0.85)
    {output buffers};

  \node[note, text width=12.0cm] at (0,-1.75)
    {Capacity is the sum of simultaneously live objects; traffic depends on
     how often each object is read or written.  These are related but distinct
     limits.};
\end{tikzpicture}%
}
\caption{Inference stages and tensor lifetimes for one autoregressive
request.  The diagram separates persistent, growing and transient state; it
does not imply that every runtime materialises each buffer separately.
Author's synthesis, informed by Pope et al., FlexGen and
PagedAttention~\cite{pope2023efficient,sheng2023flexgen,kwon2023pagedattention}.}
\label{fig:inference-memory-lifetimes}
\end{figure}
\FloatBarrier

Three latency and throughput concepts follow directly.  \emph{Time to first
token} (TTFT) covers the request boundary through the first generated token and
is therefore driven largely, but not exclusively, by prefill.  \emph{Inter-token
latency} (ITL), also called time per output token, measures successive decode
steps.  Throughput in tokens per second depends on whether tokens are counted
for one sequence or across a batch.  These metrics answer different questions
and must be reported with prompt length, generated length and concurrency.

\subsection{Where memory can limit CPU inference}
\label{sec:inference-pipeline-bottlenecks}

The pipeline exposes several different memory problems; referring simply to
``the memory bottleneck'' hides the mechanism.  \tabref{tab:inference-pipeline}
separates them.  Na et al.'s CPU characterisation confirms that prefill and
decode stress the machine differently~\cite{na2024llmcpu}, while analytical
work across inference hardware explains the low arithmetic intensity of
small-batch generation~\cite{pope2023efficient,yuan2024roofline}.  Evidence
from accelerator serving is included where it identifies a tensor lifetime or
allocation problem, not as a prediction of a Xeon Max speedup.

\begin{table}[htbp]
\centering
\small
\tabsetup
\caption{Memory pressure along an LLM inference request.  ``Possible limit''
names a mechanism to test, not an assumption that it dominates every model or
runtime.}
\label{tab:inference-pipeline}
\begin{tabularx}{\textwidth}{@{}P{2.15cm}P{3.05cm}YP{3.0cm}@{}}
\toprule
\headrow
\textbf{Stage} & \textbf{Principal objects} & \textbf{Possible memory limit}
& \textbf{Treatment here} \\
\midrule
Load and initialise
& model file, weights, packed copies
& Storage I/O, page faults and temporary duplication determine startup and
  whether the resident model fits a tier.
& Measured separately; excluded from steady-state token rates. \\
\addlinespace[2pt]
Tokenize and prepare
& input text, token IDs, request metadata
& Usually small in bytes, but serial runtime work can enter TTFT.
& Kept inside end-to-end TTFT; not attributed to HBM without evidence. \\
\addlinespace[2pt]
Prefill
& weights, prompt activations, KV writes
& Long prompts expose matrix--matrix parallelism and temporary storage, while
  weights are reused across positions and the initial KV cache is written.
& Prompt-evaluation time and throughput reported separately. \\
\addlinespace[2pt]
Decode linear layers
& weights, small activations
& At low batch, little reuse is available before the next layer and the
  weight set may be streamed again for every token.
& Primary candidate for a bandwidth benefit; tested against Roofline and
  counters. \\
\addlinespace[2pt]
Decode attention
& growing KV cache, new key/value
& Dense attention reads an increasing history and appends new state; capacity
  grows with context and batch, and allocation can fragment.
& Context and concurrency sweep plus resident-set and placement checks. \\
\addlinespace[2pt]
Select and detokenize
& vocabulary scores, sampling and text buffers
& Small state but serial execution, synchronisation or runtime overhead can
  cap latency even when bandwidth is available.
& Included in user-visible latency; isolated before assigning causality. \\
\bottomrule
\end{tabularx}
\end{table}
\FloatBarrier

These limits interact.  Quantization shrinks the persistent weight stream but
does not necessarily shrink the KV cache; grouped-query or multi-query
attention reduces KV state; batching reuses weights but increases request
state; and long context increases KV traffic even when the weight footprint is
unchanged.  A working set may be \emph{capacity-bound} because it does not fit,
\emph{bandwidth-bound} because useful work waits for a sustained byte rate, or
\emph{latency/overhead-bound} because too few accesses are in flight to reach
that rate.  Only the second condition predicts a proportional benefit from a
faster memory tier.  \secref{sec:inference-model} develops the quantitative
tests for distinguishing these cases after the next two sections establish
the cluster and processor on which they occur.
